% Transpiler optimization levels: qiskit 1.2.4 (nov2025 codebase) vs 2.6.0
% (nov2026 codebase), for the fvm nozzle HHL circuit on the ideal Aer simulator.
% Needs no packages beyond the LaTeX core; \input it into the body. In a
% two-column class, change table to table*.
%
% Sources (q8020-cfd-experiments):
%   Table 1: aux/fvm_euler_1d_nozzle/optlevel-hhl-nelem5-2026-11/qiskit-passes/
%            passes_qiskit1.2.4.txt and passes_qiskit2.6.0.txt (dump_passes.py)
%   Table 2: aux/fvm_euler_1d_nozzle/optlevel-hhl-nelem5-2026-11/
%            compare-nov2025-nov2026/data/per_case.csv (step 0, shots 0)
%   Runs:    results/fvm_euler_1d_solver/2026-11-07-hhl-nelem5-optlevel/
%   Gate mix: results/fvm_euler_1d_solver/nov2025-rerun/nov_pkgs/2026-11-03/
%            _9f0fbf67/719a35aa/circuit_transpile_statevector_nelem5_iter0.qpy

\begin{table}[htbp]
\centering
\caption{Passes run at each preset optimization level by
\texttt{transpile(circ, AerSimulator())}, qiskit 1.2.4 against 2.6.0. A
bullet marks a pass or behaviour that the level uses. Layout and routing are
empty at every level in both versions, because the simulator target has no
coupling map; it accepts 98 operations natively. Bookkeeping passes are left
out.}
\label{tab:qiskit-optlevel-passes}
\footnotesize
\begin{tabular}{l|cccc|cccc}
\hline
 & \multicolumn{4}{c|}{qiskit 1.2.4} & \multicolumn{4}{c}{qiskit 2.6.0} \\
Level & 0 & 1 & 2 & 3 & 0 & 1 & 2 & 3 \\
\hline
\multicolumn{9}{l}{\textit{Synthesis and translation}} \\
\texttt{HighLevelSynthesis}, \texttt{BasisTranslator}
  & $\bullet$ & $\bullet$ & $\bullet$ & $\bullet$
  & $\bullet$ & $\bullet$ & $\bullet$ & $\bullet$ \\
\texttt{UnitarySynthesis}
  & $\bullet$ & $\bullet$ & $\bullet$ & $\bullet$
  & $\bullet$ & $\bullet$ & $\bullet$ & $\bullet$ \\
\hline
\multicolumn{9}{l}{\textit{Removal and cancellation}} \\
\texttt{ElidePermutations}
  & & & $\bullet$ & $\bullet$ & & & $\bullet$ & $\bullet$ \\
\texttt{RemoveDiagonalGatesBeforeMeasure}
  & & & $\bullet$ & $\bullet$ & & & $\bullet$ & $\bullet$ \\
\texttt{RemoveIdentityEquivalent}
  & & & & & & & $\bullet$ & $\bullet$ \\
\texttt{InverseCancellation}
  & & $\bullet$ & $\bullet$ & $\bullet$ & & $\bullet$ & $\bullet$ & $\bullet$ \\
\texttt{CommutativeCancellation}
  & & & $\bullet$ & $\bullet$ & & & $\bullet$ & $\bullet$ \\
\texttt{Optimize1qGatesDecomposition}
  & & $\bullet$ & $\bullet$ & $\bullet$ & & $\bullet$ & $\bullet$ & $\bullet$ \\
\hline
\multicolumn{9}{l}{\textit{Two-qubit blocks}} \\
\texttt{Collect2qBlocks}
  & & & $\bullet$ & $\bullet$ & & & & \\
\texttt{ConsolidateBlocks}$^{a}$
  & & & $\bullet$ & $\bullet$ & & & $\bullet$ & $\bullet$ \\
\texttt{Split2QUnitaries}
  & & & $\bullet$ & $\bullet$ & & & $\bullet$ & $\bullet$ \\
\texttt{TwoQubitPeepholeOptimization}
  & & & & & & & $\bullet$ & $\bullet$ \\
\hline
\multicolumn{9}{l}{\textit{Control}} \\
\texttt{ContractIdleWiresInControlFlow}$^{b}$
  & & & & & & $\bullet$ & $\bullet$ & $\bullet$ \\
Repeat until no change
  & & $\bullet$ & $\bullet$ & & & $\bullet$ & $\bullet$ & \\
Keep the best circuit found
  & & & & $\bullet$ & & & & $\bullet$ \\
Default when no level is given
  & & $\bullet$ & & & & & $\bullet$ & \\
\hline
\end{tabular}

\smallskip
\begin{minipage}{0.92\linewidth}
\footnotesize\raggedright
$^{a}$In 2.6.0, \texttt{ConsolidateBlocks} collects its own blocks, and
\texttt{TwoQubitPeepholeOptimization} takes over the two-qubit block
resynthesis that 1.2.4 does with \texttt{Collect2qBlocks},
\texttt{ConsolidateBlocks} and \texttt{UnitarySynthesis}.
$^{b}$Acts only on control-flow operations; the HHL circuit has none.
\end{minipage}
\end{table}

\begin{table}[htbp]
\centering
\caption{Transpiled HHL circuit for the 1D nozzle at 5 elements, Newton step 0
(14 qubits, exact statevector run), by optimization level. The gate count
includes all gate types. Transpile time is the mean over the 8 Newton steps
of the run, in s, with 8 runs sharing a 12-core laptop.}
\label{tab:qiskit-optlevel-hhl}
\footnotesize
\begin{tabular}{c|rrr|rrr|cc}
\hline
 & \multicolumn{3}{c|}{Depth} & \multicolumn{3}{c|}{Gates}
 & \multicolumn{2}{c}{Transpile (s/step)} \\
Level & 1.2.4 & 2.6.0 & change & 1.2.4 & 2.6.0 & change & 1.2.4 & 2.6.0 \\
\hline
0 & 689\,284 & 702\,509 & $+1.9\%$ & 689\,633 & 734\,779 & $+6.5\%$
  & 19.6 & 0.68 \\
1$^{*}$ & 689\,026 & 638\,248 & $-7.4\%$ & 689\,112 & 649\,091 & $-5.8\%$
  & 17.6 & 1.5 \\
2$^{\dagger}$ & 689\,026 & 596\,687 & $-13.4\%$ & 689\,112 & 641\,137
  & $-7.0\%$ & 53.6 & 24.3 \\
3 & 689\,026 & 596\,687 & $-13.4\%$ & 689\,112 & 641\,137 & $-7.0\%$
  & 50.8 & 56.5 \\
\hline
\end{tabular}

\smallskip
\begin{minipage}{0.92\linewidth}
\footnotesize\raggedright
$^{*}$Default level on qiskit 1.2.4. $^{\dagger}$Default level on qiskit
2.6.0. A replicate run gives transpile times up to about 1.5 times these. The
two codebases also use different ports of the HHL circuit library (jw676
\texttt{a94cf63} and Q8020 \texttt{99231b85}), so the table compares software
stacks; the generated circuit has the same structure in both (5 layers, 19
instructions).
\end{minipage}
\end{table}

In both qiskit versions the four optimization levels share one structure
(Table~\ref{tab:qiskit-optlevel-passes}). Level~0 only synthesizes the
high-level HHL blocks (phase estimation, the reciprocal rotation and inverse
phase estimation) and translates them into gates the target accepts. Level~1
adds cancellation of inverse gate pairs and a loop that merges runs of
single-qubit gates. Levels~2 and~3 add an init stage that also elides
permutations, removes diagonal gates before measurement, applies
commutation-aware cancellation and consolidates two-qubit blocks. They also
resynthesize two-qubit blocks, once at level~2 and inside the loop at
level~3. Level~3 then keeps the best circuit it has seen rather than the first
fixed point. A simulator target needs no layout or routing, so the levels
differ only in these passes.

Three things change from qiskit 1.2.4 to 2.6.0. First, the default level rises
from 1 to 2. Second, a single \texttt{TwoQubitPeepholeOptimization} pass
replaces the collect, consolidate and resynthesize sequence. It keeps a
resynthesized block only when the block has fewer two-qubit gates, then
higher estimated fidelity, then fewer gates. Third, levels~2 and~3 gain
\texttt{RemoveIdentityEquivalent}, which drops gates within tolerance of the
identity.

On qiskit 1.2.4, levels~1 to~3 return the same circuit, and level~0 is only
258 layers deeper (Table~\ref{tab:qiskit-optlevel-hhl}). At the default level,
99.8\% of the 689\,126 gates in the step-0 circuit are \texttt{cu},
\texttt{ccx} and \texttt{cp} (346\,800, 226\,696 and 114\,296), which the
simulator executes natively, and only 679 are single-qubit gates. This fits
the higher levels finding nothing to remove: there is little single-qubit work
to merge, and the three-qubit \texttt{ccx} gates break up the two-qubit
blocks. On qiskit 2.6.0, each level up to~2 shortens the circuit: the depth
goes from 702\,509 to 638\,248 to 596\,687. Level~3 gives the same step-0
circuit as level~2 at 2.3 times the transpile time. At each version's default
level, the depth falls by 13.4\%, from 689\,026 to 596\,687. Version~2.6.0 is
also 11 to 29 times faster at levels~0 and~1, about 2 times faster at level~2
and about as fast at level~3.

On the noiseless simulator the level changes the cost, not the answer. The
state fidelity of the HHL update against the classical solve agrees to six
digits across all eight version and level combinations (mean 0.999815 over
Newton steps 0 to 6), and the seeded 2.6.0 shot runs are bit-identical across
levels. One caution applies to statevector extraction from the transpiled
circuit. At levels~2 and~3, \texttt{ElidePermutations} and
\texttt{RemoveDiagonalGatesBeforeMeasure} may change the state just before
measurement without changing measurement statistics. They changed nothing
here, but they run at the default level on 2.x and did not on 1.2.4.
